Multispectral UAV reconstruction is not only a rendering task; it is also a shared-support problem. A usable model must preserve per-band image quality while keeping G/R/RE/NIR predictions on the same 3D support so that downstream indices (e.g., NDVI, GNDVI, NDRE) remain stable. Recent multispectral Gaussian methods have shown that joint spectral rendering and vegetation-index synthesis are feasible, but UAV deployment still requires an explicit shared-support formulation under raw capture, rectification, and model-level product synthesis.

We propose UMGS, a two-stage UAV multispectral Gaussian splatting framework for spectral-index synthesis. Stage 1 learns an RGB anchor with standard Gaussian splatting. Stage 2 transfers spectral appearance while freezing geometry, opacity, and Gaussian identity. We study two locked variants: UMGS-I, independent per-band transfer from the shared RGB anchor, and UMGS-J, joint band-bank transfer in one unified multispectral state with a single authoritative checkpoint and per-band exported evaluation views.

We evaluate on a 10-scene UAV-focused benchmark, including 7 raw UAV captures and 3 aligned aerial multispectral scenes. UMGS-I and UMGS-J are numerically near-tied on this benchmark and show matched relative-depth behavior under held-out probe views. On aligned aerial scenes, common-mask comparison against MS-Splatting is competitive and scene-dependent. Ablations confirm the role of geometry locking: Scratch-MS and Unfrozen-support can improve selected single-band metrics, but they break the shared Gaussian support needed for model-level multispectral products and index synthesis.

Overall, the results support geometry locking as a practical design for UAV multispectral Gaussian splatting and stable spectral-index synthesis.
